% Generated draft: requires mathematical and source review.
\documentclass[11pt]{article}
\usepackage[margin=25mm]{geometry}
\usepackage{fontspec}
\setmainfont{Latin Modern Roman}
\usepackage{amsmath,amssymb,mathtools,unicode-math}
\setmathfont{Latin Modern Math}
\usepackage{enumitem,needspace,xurl,hyperref}
\setlength{\parindent}{0pt}
\setlength{\parskip}{5pt}
\setlength{\emergencystretch}{3em}
\newcommand{\N}{\mathbb{N}}
\newcommand{\Z}{\mathbb{Z}}
\newcommand{\Q}{\mathbb{Q}}
\newcommand{\R}{\mathbb{R}}
\newcommand{\C}{\mathbb{C}}
\newcommand{\F}{\mathbb{F}}
\newcommand{\A}{\mathbb{A}}
\newcommand{\PP}{\mathbb{P}}
\newcommand{\E}{\mathbb{E}}
\newcommand{\eps}{\varepsilon}
\newcommand{\dd}{\,\mathrm{d}}
\begin{document}
\section*{UnsolvedMath ID: 30001868}
\textbf{Source problem code:} catalogue entry
\subsection{Definitions and mathematical statement}
The supplied question is: \emph{Which subalgebras of a polynomial ring have a SAGBI basis?}

To specify the usual terminology over a field, let \(k\) be a field, let \(n\geq 1\) be an integer, and let \(A\subseteq k[x_1,\ldots,x_n]\) be a unital \(k\)-subalgebra. Fix a monomial order \(\prec\), meaning a total well-order on monomials, compatible with multiplication and with \(1\) least. For a nonzero polynomial \(f\), write \(\operatorname{LM}_{\prec}(f)\) for its largest monomial under \(\prec\). Define the initial algebra by
\[
\operatorname{in}_{\prec}(A)
 = k\bigl[\operatorname{LM}_{\prec}(f):f\in A\setminus\{0\}\bigr].
\]
A subset \(B\subseteq A\setminus\{0\}\) is a \emph{SAGBI basis} of \(A\) with respect to \(\prec\) if
\[
\operatorname{in}_{\prec}(A)
 = k\bigl[\operatorname{LM}_{\prec}(b):b\in B\bigr].
\]
It is a finite SAGBI basis when \(B\) is finite.

The supplied literature assessment interprets the target as a characterization of subalgebras admitting \emph{finite} SAGBI bases. Under that interpretation, for a specified \(k\), \(n\), and \(\prec\), the characterization problem concerns the condition
\[
\exists\,B\subseteq A\setminus\{0\},\quad |B|<\infty,
\qquad
\operatorname{in}_{\prec}(A)
 =k\bigl[\operatorname{LM}_{\prec}(b):b\in B\bigr].
\]
If the order is instead allowed to vary with \(A\), the condition becomes
\[
\exists\,\prec\ \exists\,B\subseteq A\setminus\{0\},\quad |B|<\infty,
\qquad
\operatorname{in}_{\prec}(A)
 =k\bigl[\operatorname{LM}_{\prec}(b):b\in B\bigr].
\]
The record does not determine which order quantifier is intended. It also does not specify the coefficient domain, the number of variables, or whether the subalgebras under consideration must be finitely generated. The field-based definitions above therefore explain a conventional setting rather than supply missing source hypotheses. Moreover, the question itself omits finiteness: if arbitrary infinite SAGBI bases are permitted under the displayed definition, \(A\setminus\{0\}\) is such a basis by construction. Thus the finite interpretation comes from the supplied assessment and requires confirmation against the source context.
\subsection{Short English statement}
Characterize polynomial subalgebras whose initial algebras are generated by the leading monomials of finitely many of their elements. Finiteness is indicated by the supplied literature assessment but is absent from the quoted question. The coefficient domain, restrictions on the subalgebras, and whether the monomial order is fixed or may be chosen remain unresolved.
\subsection{Sources}
\begin{itemize}
\item Local UnsolvedMath registry, supplied record with permanent ID \(30001868\), title \emph{Polynomial Subalgebras Admitting SAGBI Bases}, exact source problem code \texttt{catalogue entry}, in the Oberwolfach Reports Open Problems collection. The supplied proposed year is 2011; no individual proposer is named.
\item Supplied source citation: \emph{Mini-Workshop: New Developments in Newton-Okounkov Bodies (2012)}, Oberwolfach Reports, pp.\ 2327--2363, DOI \texttt{10.4171/owr/2011/41}. Original OWR report: \url{https://doi.org/10.4171/owr/2011/41}. The background describes Oberwolfach Reports as records of mathematical workshops and their open-problem discussions. The citation's parenthetical year 2012 differs from the supplied proposed year 2011 and the year component of the DOI; this discrepancy is retained without resolution.
\item Supplied literature source, labeled ``SAGBI bases survey'': \url{https://www.math.ubc.ca/~reichst/sagbi.pdf}. No author, publication date, or fuller bibliographic citation is supplied for this link.
\item Supplied literature assessment dated \texttt{2026-08-21}: the record labels the problem open and describes characterization of subalgebras admitting finite SAGBI bases as a recognized open problem, while stating that algorithms and important classes are known. This is supplied dated curation, not an independently verified literature assessment.
\item The record labels the statement ``exact'' and reports that its mathematical content was checked against the official Oberwolfach report, with notation, context, and question form normalized. These are supplied verification claims; no independent source inspection is asserted here.
\end{itemize}
\end{document}
